% Figura Anexo I: Mundell-Fleming con tipo de cambio flexible y movilidad perfecta: (a) política monetaria; (b) política fiscal
\documentclass[tikz,border=3pt]{standalone}
\usepackage[T1]{fontenc}
\usepackage[utf8]{inputenc}
\usetikzlibrary{arrows.meta,calc}
\begin{document}
\begin{tikzpicture}[>={Latex[length=2.2mm]}, line cap=round, font=\footnotesize,
  eje/.style={->, thick},
  curva/.style={very thick},
  nueva/.style={very thick, blue!55!black},
  guia/.style={densely dashed, gray!80},
  pto/.style={circle, fill=black, inner sep=1.1pt}]
  % ---------- Panel (a): política monetaria ----------
  \node[font=\footnotesize\bfseries] at (3.6,5.15) {(a) Política monetaria, tipo flexible};
  \draw[eje] (0,0) -- (7.3,0) node[below left=1pt]{$Y$};
  \draw[eje] (0,0) -- (0,4.9) node[below right=2pt]{$r$};
  \draw[curva, gray!45!black] (0.3,2.2) -- (6.7,2.2) node[right]{$BP\ (r=r^{*})$};
  \draw[curva] (0.3,0.715) -- (6.3,4.02) node[right]{$LM_0$};
  \draw[nueva] (0.9,0.245) -- (6.8,3.49) node[right]{$LM_1$};
  \draw[curva] (0.5,4.075) -- (5.4,0.4) node[below right=-3pt]{$IS_0$};
  \draw[nueva] (1.6,4.34) -- (6.85,0.4025) node[below right=-3pt]{$IS_1$};
  \node[pto] (A0) at (3,2.2) {};
  \node[above left=-1pt of A0, inner sep=1pt] {$E_0$};
  \node[circle, draw=black, fill=white, inner sep=1.1pt] (Ap) at (3.615,1.738) {};
  \node[below left=-2pt of Ap, inner sep=1pt] {$E'$};
  \node[pto, fill=blue!55!black] (A1) at (4.4545,2.2) {};
  \node[above right=-1pt of A1, inner sep=1pt, blue!55!black] {$E_1$};
  \draw[->, shorten >=4pt, shorten <=4pt, gray!60!black] (A0) -- (Ap) node[midway, below left=-2pt, black]{\scriptsize 1};
  \draw[->, shorten >=4pt, shorten <=4pt, gray!60!black] (Ap) -- (A1) node[midway, below right=-2pt, black]{\scriptsize 2};
  \draw[guia] (3,2.2) -- (3,0) node[below]{$Y_0$};
  \draw[guia] (4.4545,2.2) -- (4.4545,0) node[below]{$Y_1$};
  \node[align=left, anchor=north west] at (0.3,-0.75) {\scriptsize 1: $\Delta M$ baja $r$ y salen capitales $\Rightarrow$ depreciación\\[-2pt]\scriptsize 2: la depreciación eleva las exportaciones netas ($IS_0\to IS_1$)};
  % ---------- Panel (b): política fiscal ----------
  \begin{scope}[xshift=9.0cm]
    \node[font=\footnotesize\bfseries] at (3.6,5.15) {(b) Política fiscal, tipo flexible};
    \draw[eje] (0,0) -- (7.3,0) node[below left=1pt]{$Y$};
    \draw[eje] (0,0) -- (0,4.9) node[below right=2pt]{$r$};
    \draw[curva, gray!45!black] (0.3,2.2) -- (6.7,2.2) node[right]{$BP\ (r=r^{*})$};
    \draw[curva] (0.3,0.715) -- (6.3,4.02) node[right]{$LM_0$};
    \draw[curva] (0.5,4.075) -- (5.4,0.4) node[below right=-3pt]{$IS_0$};
    \draw[nueva] (1.6,4.34) -- (6.85,0.4025) node[below right=-3pt]{$IS_1$};
    \node[pto] (B0) at (3,2.2) {};
    \node[below left=-1pt of B0, inner sep=1pt] {$E_0=E_1$};
    \node[circle, draw=black, fill=white, inner sep=1.1pt] (Bp) at (3.838,2.661) {};
    \node[above right=-2pt of Bp, inner sep=1pt] {$E'$};
    \draw[->, shorten >=4pt, shorten <=4pt, gray!60!black] (B0) -- (Bp) node[midway, above left=-2pt, black]{\scriptsize 1};
    \draw[->, shorten >=4pt, shorten <=4pt, gray!60!black] (Bp) to[bend left=18] (B0);
    \node[gray!40!black] at (4.45,1.95) {\scriptsize 2};
    \draw[guia] (3,2.2) -- (3,0) node[below]{$Y_0$};
    \node[align=left, anchor=north west] at (0.3,-0.75) {\scriptsize 1: $\Delta G$ presiona $r$ al alza y entran capitales $\Rightarrow$ apreciación\\[-2pt]\scriptsize 2: la apreciación hunde las exportaciones netas: la $IS$ vuelve ($Y$ a $Y_0$)};
  \end{scope}
\end{tikzpicture}
\end{document}
